% ============================================================
% Recreational Marijuana Legalization and Drug-Overdose
% Mortality: A Cross-Border Distance Design
% Formatted to Economic Inquiry (Wiley/WEAI) submission requirements.
% REVISION (2026-07-27): full rewrite bringing this draft up to date
% with the current, most-developed version of the paper. Adds: the
% pre-trends/mean-reversion correction and its resulting change to the
% preferred specification (previously this draft still showed the raw,
% uncorrected coefficients); the county fixed-effects check; the
% period-specific-effects table (including the investigated and
% resolved 2023-2024 sign reversal); the persistence check; the
% diff-in-discontinuities appendix check; the full robustness-check
% summary table; the corrected (unified-panel) pre-trends numbers; and
% two brand-new robustness checks: an estimator comparison against
% two-way fixed effects and Callaway-Sant'Anna (2021), and a
% dispensary-level check using real geocoded license-opening dates in
% place of state-legalization-date distance.
% Compile: pdfLaTeX -> BibTeX -> pdfLaTeX -> pdfLaTeX
% ============================================================
\documentclass[11pt]{article}

\usepackage[margin=1in]{geometry}
\usepackage{setspace}
\usepackage{amsmath,amssymb}
\usepackage{graphicx}
\usepackage{booktabs}
\usepackage{caption}
\usepackage[authoryear,round]{natbib}
\usepackage[hidelinks]{hyperref}
\usepackage{enumitem}
\usepackage{titlesec}

\doublespacing
\titleformat{\section}{\normalfont\bfseries\large}{\thesection.}{0.5em}{}
\titleformat{\subsection}{\normalfont\bfseries}{\thesubsection}{0.5em}{}

\title{}
\author{}
\date{}

\begin{document}

% ============================================================
% TITLE PAGE
% ============================================================
\begin{titlepage}
\begin{center}
\vspace*{1.5cm}
{\LARGE \bfseries Recreational Marijuana Legalization and\\
Drug-Overdose Mortality: A Cross-Border\\
Distance Design\par}
\vspace{0.3cm}
{\large \itshape A Stacked Multi-Period Panel Using Every State's Legalization Date, 2020--2025\par}
\vspace{1.2cm}
{\large Author name(s) withheld for review\par}
\vspace{0.3cm}
{\normalsize Working Paper --- July 2026\par}
\end{center}

\vspace{1.2cm}
\noindent\textbf{Abstract} \\[4pt]
\noindent We test whether proximity to a legal recreational-marijuana
market reduces drug-overdose mortality, using a stacked multi-period
panel of counties in 31 never-legalizing states and every state's
actual legalization date, 2020--2025. Switcher counties see overdose
deaths fall 2.32 per 100,000 more than controls ($p=0.004$); a
continuous per-mile specification agrees in sign and is itself
significant ($p<0.001$), surviving even county fixed effects, unlike
the switcher indicator. A formal pre-trends test finds and corrects a
mean-reversion problem in the raw comparison. A medical-legalization
placebo check is null. Comparing our design against two-way fixed
effects and Callaway-Sant'Anna (2021) shows neither off-the-shelf
staggered-adoption estimator matches our stacked-cohort result exactly
-- informatively so, since Callaway-Sant'Anna's own pre-treatment
coefficients independently confirm the same mean-reversion problem we
correct for. A further check replacing state-legalization timing with
real, geocoded dispensary-opening dates for a 9-state subsample
reproduces the same sign. The effect implies roughly 755 averted
deaths and 22,500 years of life saved. \\[10pt]
\noindent\textbf{JEL Classification:} I18, I12, K42, C23 \\
\noindent\textbf{Keywords:} recreational marijuana; opioid overdose
mortality; cross-border spillover; staggered difference-in-differences;
drug policy

\vfill
\noindent\textbf{Abbreviations} \\[2pt]
\noindent
\begin{tabular}{@{}ll@{}}
ATT    & Average treatment effect on the treated \\
CCT    & Calonico-Cattaneo-Titiunik (optimal-bandwidth selector) \\
CS     & Callaway-Sant'Anna (group-time treatment-effect estimator) \\
DiD    & Difference-in-differences \\
FIPS   & Federal Information Processing Standard (county code) \\
LE     & Life expectancy \\
NCHS   & National Center for Health Statistics \\
SE     & Standard error \\
TWFE   & Two-way fixed effects \\
VSRR   & Vital Statistics Rapid Release \\
YLL    & Years of life lost/saved \\
\end{tabular}

\vspace{0.5cm}
\noindent{\footnotesize A separate Disclosure Statement accompanies this
submission per Economic Inquiry policy and is not reproduced in this file.}
\end{titlepage}

% ============================================================
% 1. INTRODUCTION
% ============================================================
\section{Introduction}

Does proximity to a legal recreational-marijuana market change
drug-overdose mortality for residents who never cross the border
themselves? The evidence on marijuana-opioid substitution is
unsettled. \citet{Bachhuber2014} found early evidence that state
medical-marijuana laws were associated with fewer opioid overdose
deaths. \citet{Shover2019} showed the association reverses sign once
later years of data are included, and \citet{PowellPaculaJacobson2018}
found the effect, where present, is specific to dispensary-based legal
designs rather than medical-marijuana laws generally.
\citet{MathurRuhm2022} conclude that legal retail dispensaries --
medical or recreational -- are, if anything, associated with
\emph{higher} opioid mortality. Most recently, \citet{Denkyirah2026}
use a state-level two-way-fixed-effects design focused specifically on
\emph{recreational} legalization and find a robust decline of
approximately 3.5 opioid overdose deaths per 100,000 -- a magnitude we
return to below.

Every paper above identifies off the \emph{timing} of state policy
adoption. None uses \emph{geography}: whether a county's proximity to
an external legal market, in a state that has not itself legalized,
predicts its own overdose trajectory. This is the design
\citet{HaoCowan2020} developed for marijuana-possession arrests. We
apply it to overdose mortality -- to our knowledge a new combination --
using both a binary switcher indicator and a continuous, per-mile
distance measure, built as a stacked multi-period panel that uses all
24 states' actual legalization dates rather than collapsing a
multi-year staggered rollout into two snapshots. An earlier,
two-snapshot version of this design found a null, sign-flipped
continuous estimate that we could not resolve with a bandwidth-based
fix; the multi-period panel resolves it directly, and further makes
possible a formal pre-trends test, a persistence (post-switch
dynamics) check, and a county fixed-effects robustness test -- all
reported below, together with a placebo-style check that swaps in
medical- for recreational-marijuana legalization timing, a comparison
against two canonical staggered-adoption estimators (two-way fixed
effects and Callaway-Sant'Anna), and a further robustness check that
replaces our county-centroid distance measure with distance to real,
geocoded, individually-dated dispensaries where such records are
publicly available.

% ============================================================
% 2. DATA AND EMPIRICAL STRATEGY
% ============================================================
\section{Data and empirical strategy}

We use 2,331 counties in the 31 U.S. states with no
recreational-marijuana retail market as of our latest
snapshot.\footnote{The companion project's original 32-state
never-legalizing panel included Maine; Maine's retail market opened
October 2020, before even our earliest snapshot, so its counties are
excluded here.} Each county's distance to the nearest legal-market
county is computed via Haversine distance between county centroids
(Census Gazetteer), with ``legal-market'' defined by first retail
sales date, for every December 2020--2025 -- the full window CDC
publishes county-level 12-month-ending overdose data for. Rather than
compare two fixed snapshots, we use all 24 states' actual
retail-opening dates and build a stacked panel of year-over-year
transitions. In transition $t-1\to t$, a county is a ``new switcher''
if its distance to the nearest legal-market county crossed from at
least 50 miles to below 50; a ``control'' if it stayed at least 50
miles away at both endpoints. Once a county switches it is dropped
from every later transition, avoiding the already-treated-unit-as-
control bias documented in the staggered difference-in-differences
literature \citep{GoodmanBacon2021}. Counties within 50 miles as of
the December-2020 baseline (116 of 2,331) are ``always-near'' and
dropped throughout.

The unit of observation is a county-transition pair $(i,t)$: county
$i$ contributes one row for each one-year transition
$t\in\{\text{2020--21},\ldots,\text{2024--25}\}$ it is ``at risk'' of
switching in -- every transition up to and including the one in which
it switches, if it ever does, after which it is dropped, as just
described. $\Delta Y_{it}$ is the change in county $i$'s overdose
death rate per 100,000 between the start and end of transition $t$.
$\text{Switcher}_{it}$ is a binary indicator equal to 1 if and only if
county $i$'s distance to the nearest legal-market county crossed from
at least 50 miles to below 50 miles \emph{during} transition $t$, so a
given county contributes at most one $\text{Switcher}_{it}=1$
observation across the whole panel; it equals 0 for every transition a
county has not yet switched in, including every transition for
counties that never switch. $\Delta\text{Dist}_{it}$ is the raw change
in miles to the nearest legal-market county over transition $t$,
defined for every county-transition row regardless of switcher status
(control counties can still move closer or farther without crossing
the threshold). $\delta_t$ denotes transition (period) fixed effects,
common to all counties observed in a given transition. Standard errors
are clustered by state, the level at which most of the underlying
legal-market-timing variation occurs. We pool all (county, transition)
observations into one stacked regression
\citep{CengizDubeLindnerZipperer2019}:
\begin{equation}
  \Delta Y_{it} = \alpha + \gamma \cdot \text{Switcher}_{it} + \delta_t + \varepsilon_{it}, \qquad
  \Delta Y_{it} = \alpha + \beta \cdot \Delta\text{Dist}_{it} + \delta_t + \varepsilon_{it}.
  \label{eq:main}
\end{equation}
Equation~\ref{eq:main} imposes one switcher effect $\gamma$ and one
distance effect $\beta$ common to all five transitions. We also
estimate a more general version that lets both vary freely by period,
via a full interaction with the transition dummies:
\begin{equation}
  \Delta Y_{it} = \alpha + \sum_t \gamma_t \cdot (\text{Switcher}_{it} \times \mathbb{1}[t])
                        + \sum_t \beta_t \cdot (\Delta\text{Dist}_{it} \times \mathbb{1}[t])
                        + \delta_t + \varepsilon_{it},
  \label{eq:period}
\end{equation}
which nests Equation~\ref{eq:main} as the special case
$\gamma_t=\gamma$ and $\beta_t=\beta$ for all $t$;
Section~\ref{sec:periodeffects} reports the period-specific
coefficients and tests whether pooling is appropriate.

As of December 2020, ten states had an open retail market; by December
2025, 23 did (Figure~\ref{fig:rollout}). Virginia has never opened a
general retail market and so never enters the legal-market set.
County-level overdose mortality comes from CDC's Vital Statistics
Rapid Release (VSRR) provisional county data (data.cdc.gov); CDC
WONDER is not reachable from our environment. NCHS confidentiality
rules suppress county-transition cells with 1--9 deaths at either
endpoint, leaving 4,259 of 10,313 potential observations (177 switch
events).

As a placebo-style check, we rebuild the identical stacked design
using each state's medical-marijuana access date (first regulated
medical dispensary sales, compiled for 39 states) in place of the
recreational-retail date, restricting the county universe to the 17 of
31 states not yet medically legal as of the December-2020 baseline.

\begin{figure}[htbp]
  \centering
  \includegraphics[width=\textwidth]{figures/fig1_rollout_maps_multiperiod.png}
  \caption{Legal-market rollout (by first retail sales date) and 50-mile treatment frontier, 2020--2025.}
  \label{fig:rollout}
\end{figure}

% ============================================================
% 3. RESULTS
% ============================================================
\section{Results}

\begin{table}[htbp]
\centering
\caption{Effect on drug-overdose mortality, preferred specification}
\label{tab:main}
\begin{tabular}{@{}lcc@{}}
\toprule
 & Switcher (binary) & Distance, per mile \\
\midrule
Overdose rate per 100k & $-2.320^{***}$ (0.800) & $+0.0122^{***}$ (0.0034) \\
 & $p=0.004$ & $p<0.001$ \\
\bottomrule
\end{tabular}
\vspace{2pt}
\begin{minipage}{0.85\textwidth}
\footnotesize\emph{Notes:} Outcome is the transition-level change in
overdose deaths per 100{,}000. Both columns include transition fixed
effects and a control for each county's baseline (2020) overdose-rate
level (Section~\ref{sec:pretrends}). SE clustered by state. $N=3{,}794$.
The distance coefficient is signed as the effect of an
\emph{increase} in distance, so it agrees in direction with the
(negative) switcher coefficient. $^{***}p<0.01$.
\end{minipage}
\end{table}

\begin{figure}[htbp]
  \centering
  \includegraphics[width=0.85\textwidth]{figures/fig2_overdose_coefplot_multiperiod.png}
  \caption{Effect on overdose mortality, recreational (primary) vs. medical (placebo) design (95\% CI).}
  \label{fig:coef}
\end{figure}

Switcher counties see overdose deaths fall 2.32 more per 100,000 than
controls in the same transition ($p=0.004$; Table~\ref{tab:main}) --
somewhat below \citet{Denkyirah2026}'s state-level estimate of roughly
3.5, though the same direction and order of magnitude. The continuous
per-mile specification agrees in direction and is itself significant
($p<0.001$).

\subsection{Pre-trends check and robustness to fixed effects}
\label{sec:pretrends}

Because eventual-switcher counties contribute observations to every
transition before they switch, we can test whether their overdose
trend was already diverging from never-switching counties
\emph{before} switching. The raw version of this test fails: leads on
event-time indicators (one and two-plus transitions before switching)
are large and jointly significant ($p<0.001$), comparable in magnitude
to the switch-transition effect itself, even though transition fixed
effects are already present. We traced this to mean reversion --
counties with unusually high 2020 overdose rates mechanically decline
more -- and adding each county's baseline (2020) rate as a control
resolves it (joint F-test on both leads, $p=0.318$; Figure~\ref{fig:pretrends}).
We also checked Census-division-by-transition fixed effects (a
geographic-clustering explanation); this also passes the joint test on
this panel but leaves the switch coefficient less precisely estimated
and proved sensitive to exact sample construction across our
diagnostic checks, so we prefer the baseline-rate control
(Table~\ref{tab:main}).

\begin{figure}[htbp]
  \centering
  \includegraphics[width=0.85\textwidth]{figures/fig5_eventstudy_pretrends.png}
  \caption{Pre-trends check: raw vs. mean-reversion-adjusted event-study coefficients.}
  \label{fig:pretrends}
\end{figure}

\begin{table}[htbp]
\centering
\caption{Robustness of main coefficients to fixed-effects controls}
\label{tab:fe}
\begin{tabular}{@{}lcccc@{}}
\toprule
 & \multicolumn{2}{c}{Switcher (binary)} & \multicolumn{2}{c}{Distance (per mile)} \\
Specification & Coef. & $p$ & Coef. & $p$ \\
\midrule
No controls & $-3.151$ & 0.005 & $+0.0128$ & $<0.001$ \\
+ Baseline overdose-rate level (preferred) & $-2.320^{***}$ & 0.004 & $+0.0122$ & $<0.001$ \\
+ State fixed effects & $-1.411$ & 0.144 & $+0.0129$ & $<0.001$ \\
+ County fixed effects & $-0.999$ & 0.532 & $+0.0122^{**}$ & 0.023 \\
\bottomrule
\end{tabular}
\vspace{2pt}
\begin{minipage}{0.92\textwidth}
\footnotesize\emph{Notes:} Each row is a separate regression on the
same $N=3{,}794$ sample; SE clustered by state. County fixed effects
use only within-county variation across a given county's (at most
five) transition-rows, the most demanding test we run -- 225 of 1,249
counties in this sample have only one transition-row and so contribute
no within-county variation. $^{**}p<0.05$, $^{***}p<0.01$.
\end{minipage}
\end{table}

Table~\ref{tab:fe} shows the two specifications respond very
differently to fixed-effects controls. The continuous per-mile
coefficient is essentially unmoved by state \emph{or} county fixed
effects and remains significant throughout. The binary switcher
coefficient survives the baseline-rate control but loses significance
under state fixed effects and further under county fixed effects,
indicating a meaningful share of its identifying variation comes from
\emph{which} counties or states happen to border a soon-to-open
market, not from finer within-unit proximity differences. We therefore
treat the continuous per-mile result as this paper's primary
quantitative claim, with the switcher result as corroborating but
comparatively more fragile.

\subsection{Estimator robustness: two-way fixed effects and Callaway-Sant'Anna}
\label{sec:twfecs}

The stacked-cohort estimator in Equation~\ref{eq:main} is one
particular way to handle staggered treatment timing; it is not the
only one, and the staggered difference-in-differences literature has
shown that the more familiar two-way fixed-effects (TWFE) regression
can be badly biased under staggered adoption with heterogeneous
treatment effects \citep{GoodmanBacon2021}. We check whether our
result is an artifact of the stacking approach itself by re-estimating
it two more ways on the same underlying treatment definition (state
legal-market distance, county crosses below 50 miles): a canonical
TWFE regression, and the doubly-robust group-time estimator of
\citet{CallawaySantAnna2021}.

Both require a genuine county-by-year panel with an \emph{absorbing}
treatment indicator (once a county crosses the 50-mile threshold it
stays ``treated'' in every subsequent year, rather than being dropped,
as the stacked design does) and, unlike the stacked design's
transition-by-transition suppression handling, a county with any
suppressed year cannot contribute at all. This is a materially
stricter data requirement: of the same 2,215 never-legal, not-always-near
counties used elsewhere, only 614 have non-suppressed overdose data in
every one of the six panel years (2020--2025), giving a balanced panel
of 3,684 county-year observations with three identified treatment
cohorts (39 counties first crossing in 2022, 41 in 2023, 65 in 2024;
counts differ slightly from the switch counts elsewhere because this
panel requires complete data across \emph{all six} years, not just the
two years spanning a given transition).

The TWFE regression is $Y_{it} = \alpha_i + \lambda_t + \beta \cdot
\text{Treated}_{it} + \varepsilon_{it}$, county and year fixed effects,
clustered by state. The \citet{CallawaySantAnna2021} estimator computes
group-time average treatment effects $ATT(g,t)$ for each treatment
cohort $g$ against the never-treated counties, then aggregates first to
an overall average effect and separately to an event-study (dynamic)
profile by time since treatment, using the multiplier bootstrap for
inference.

\begin{table}[htbp]
\centering
\caption{Estimator comparison: TWFE and Callaway-Sant'Anna vs. the stacked-cohort design}
\label{tab:twfecs}
\begin{tabular}{@{}lccc@{}}
\toprule
Estimator & Coef. & SE & $p$ \\
\midrule
Stacked-cohort switcher (Table~\ref{tab:main}, preferred spec) & $-2.320$ & $0.800$ & $0.004$ \\
Two-way fixed effects (county + year) & $-5.727$ & $2.338$ & $0.014$ \\
Callaway-Sant'Anna (2021), overall ATT, never-treated controls & $-1.905$ & $1.929$ & $0.323$ \\
\bottomrule
\end{tabular}
\vspace{2pt}
\begin{minipage}{0.92\textwidth}
\footnotesize\emph{Notes:} SE clustered by state. TWFE and
Callaway-Sant'Anna use the balanced 614-county panel described above
($n=3{,}684$ county-years); the stacked-cohort row is reproduced from
Table~\ref{tab:main} on its own, larger sample for reference and is not
directly comparable observation-for-observation. All three use the
same underlying treatment definition (distance to nearest legal-market
county crossing 50 miles).
\end{minipage}
\end{table}

Table~\ref{tab:twfecs} shows the three estimators do not agree closely
with each other, and we think that disagreement is itself informative.
TWFE returns a coefficient more than double the stacked design's
preferred estimate ($-5.73$ vs.\ $-2.32$) -- consistent with the
well-documented direction of TWFE bias under staggered adoption, where
already-treated units act as ``forbidden comparisons'' for later-treated
units and can inflate the estimated effect. Callaway-Sant'Anna, by
contrast, returns a smaller, statistically insignificant overall ATT
($-1.91$, $p=0.32$) on the smaller balanced sample. Examining
Callaway-Sant'Anna's own event-study decomposition
(Figure~\ref{fig:cs_eventstudy}) shows why: the pre-treatment
coefficients at $t=-1$ and $t=-2$ are themselves significant
($p=0.031$ and $p<0.001$), the same mean-reversion-driven pre-trend
contamination diagnosed in Section~\ref{sec:pretrends} -- except this
estimator, run in its default form, has no analog of our baseline-rate
covariate to correct for it. Rather than treating either off-the-shelf
estimator as a referee on the stacked design, we read this as
corroborating the pre-trends finding from an entirely independent
estimator and confirming why we build a stacked cohort design with an
explicit mean-reversion control rather than applying either canonical
staggered-adoption estimator directly.

\begin{figure}[htbp]
  \centering
  \includegraphics[width=0.85\textwidth]{figures/figH_cs_eventstudy.png}
  \caption{Callaway-Sant'Anna event-study coefficients by time since treatment (never-treated control group, no covariates).}
  \label{fig:cs_eventstudy}
\end{figure}

\subsection{Robustness: real dispensary-opening dates in place of state legalization dates}
\label{sec:dispensary}

Every result so far defines ``legal-market distance'' at the state
level: a county's distance to the nearest county in a state whose
retail market is open, as of that state's own first-retail-sales date.
This is a coarse proxy for the thing that should actually matter to a
nearby resident -- physical distance to a place that sells recreational
marijuana. If the state-level measure is systematically off (e.g.\ a
county classified as ``near'' the day its neighboring state legalizes,
even though no dispensary within 50 miles has actually opened yet), it
could bias the estimated speed, though not necessarily the existence,
of the effect. We test this directly by rebuilding the identical
stacked switcher/continuous design with distance measured to the
nearest \emph{individual, currently operating} dispensary rather than
to the nearest legal-market county.

\emph{Data.} We collected official dispensary-license rosters directly
from state cannabis-regulatory agencies (not third-party aggregators)
for every state where a usable, dated, machine-readable roster is
publicly available: Ohio, New York, Massachusetts, Oregon, Illinois,
Maine, and California carry a genuine per-dispensary license-issuance
or first-retail-sale date; Washington and Colorado's own license files
carry only a renewal/expiration date, not an opening date, but both
states legalized recreational retail in 2014 -- years before our
2020--2025 panel begins -- so their dispensaries are coded as already
open throughout the panel rather than assigned a misleading proxy
date. This covers 9 of the 24 legal states; the remaining 15 either
have no public per-license date field (confirmed for Arizona,
Connecticut, Missouri, Montana, New Jersey, Delaware, Nevada, and
Vermont) or were not pursued further given time constraints (Arizona's
own market is small; several others are mid-sized). Records were
geocoded to latitude/longitude via the Census Bureau's public batch
geocoder (98.3\% match rate combining an exact street-address match for
91.7\% of records and a ZIP-centroid fallback for the remainder;
California's own regulator API already returns geocoded coordinates).
After dropping records with no usable coordinate or opening date, 5,178
of 5,825 collected dispensary licenses enter the distance calculation.

This is a genuine limitation, not a footnote to skip past: a
never-legal county whose true nearest dispensary sits in one of the 15
uncovered states is measured here as farther from a dispensary than it
actually is, which pushes that county toward the ``control'' group even
if it should be a treated ``switcher.'' This attenuates the estimated
effect and shrinks the usable switcher count rather than biasing the
sign, but it means this check should be read as confirming
\emph{direction}, not as a like-for-like replication of the main
design's precision.

\emph{Design.} Everything else is identical to Section~2: the same
2,331-county never-legal universe, the same 50-mile switcher/control
threshold, the same drop-after-switching stacked-cohort construction,
and the same CDC VSRR overdose outcome. The only change is that a
county's distance in December of year $y$ is now computed to the
nearest dispensary with an issuance date on or before December 31 of
year $y$ (Washington and Colorado dispensaries are treated as
open in every year). Fifty-three counties are ``always-near'' a real
dispensary as of the 2020 baseline under this measure -- notably fewer
than the 116 dropped under the state-level measure, since a specific
dispensary address is on average farther away than the nearest county
in \emph{any} legal state. The resulting stacked panel has 240 raw
switch events; after the same NCHS small-cell suppression rule, 84
switch events and $N=3{,}941$ county-transition observations remain.

Table~\ref{tab:dispensary} reports the continuous distance coefficient
-- our primary measure throughout this paper -- across the same
progression of fixed-effects controls as Table~\ref{tab:fe}, plus two
additional interacted specifications (state-by-transition and
census-division-by-transition fixed effects) that this smaller sample
can still support.\footnote{We could not extend this fixed-effects
menu to census-tract or calendar-season fixed effects, both requested
as robustness dimensions. Neither is estimable with the data actually
available: CDC's overdose-mortality series (VSRR and, at finer detail,
NCHS restricted files) is published only at the county level -- no
public source reports drug-overdose deaths by census tract, for
confidentiality reasons consistent with NCHS's own county-level
small-cell suppression rule already binding in this design. VSRR does
report monthly, but each monthly value is itself a 12-month
\emph{rolling} total (a moving average constructed specifically to
smooth month-to-month noise in provisional data), so it contains no
genuine within-year seasonal signal to identify a season effect
against -- a season dummy fit to a 12-month rolling series is fitting
noise, not seasonality. We use the richest fixed-effects menu the data
can actually support: county, state, census division, and
time (transition-period) fixed effects, and interactions among them.}

\begin{table}[htbp]
\centering
\caption{Dispensary-level robustness check: distance to nearest operating dispensary}
\label{tab:dispensary}
\begin{tabular}{@{}lcccc@{}}
\toprule
 & \multicolumn{2}{c}{Distance (per mile)} & \multicolumn{2}{c}{Switcher (binary)} \\
Specification & Coef. & $p$ & Coef. & $p$ \\
\midrule
No controls (transition FE only) & $+0.0230$ & 0.068 & $-5.040^{***}$ & $<0.001$ \\
+ Baseline (2020) overdose rate & $+0.0173$ & 0.082 & $-2.715^{**}$ & 0.032 \\
+ State FE & $+0.0183$ & 0.099 & $-1.936$ & 0.178 \\
+ Census division $\times$ transition FE & $+0.0171$ & 0.108 & $-2.034^{**}$ & 0.019 \\
+ State $\times$ transition FE & $+0.0051$ & 0.360 & $+3.632$ & 0.183 \\
+ County FE & $+0.0216$ & 0.124 & $-2.041$ & 0.275 \\
\bottomrule
\end{tabular}
\vspace{2pt}
\begin{minipage}{0.95\textwidth}
\footnotesize\emph{Notes:} $N=3{,}941$ throughout; SE clustered by
state. All specifications include the baseline-rate control except the
first row. Distance is signed identically to Table~\ref{tab:fe} (effect
of an \emph{increase} in distance to the nearest operating dispensary).
$^{**}p<0.05$, $^{***}p<0.01$.
\end{minipage}
\end{table}

\begin{figure}[htbp]
  \centering
  \includegraphics[width=0.85\textwidth]{figures/figG_dispensary_vs_main_coefplot.png}
  \caption{Continuous distance coefficient: dispensary-opening-date design vs. main (state-legalization-date) design, by fixed-effects specification.}
  \label{fig:dispensary}
\end{figure}

The continuous coefficient stays positive -- the same sign as the main
design throughout -- and of a broadly similar magnitude (+0.005 to
+0.023 per mile, against the main design's +0.012 to +0.013;
Figure~\ref{fig:dispensary}), and the switcher coefficient reproduces
the main design's own qualitative finding that the binary measure is
the more fragile of the two: it survives the baseline-rate control and
division-by-transition fixed effects, but loses significance under
state fixed effects, county fixed effects, and (with the ``wrong''
sign, though far from significant) state-by-transition fixed effects.
None of the continuous coefficients here reach conventional
significance at better than the 10\% level, and state-by-transition
fixed effects null it out entirely. Given how thin this specification
spreads the data -- 31 states $\times$ 5 transitions, many cells
containing only one or two counties -- we read the state-by-transition
result as a low-power, likely overfit specification rather than a
genuine contradiction, but we report it rather than omit it.

On balance, this check supports the main design's qualitative
conclusion (proximity to an actual, operating retail market is
associated with lower overdose mortality) using an entirely
independent, arguably more accurate treatment-timing measure, while
underscoring that the main design's stronger statistical significance
partly reflects its full 24-state national coverage; a 9-state,
partial-coverage version of the same underlying idea is directionally
consistent but not independently conclusive on its own.

\subsection{Period-specific effects}
\label{sec:periodeffects}

Equation~\ref{eq:main} assumes the switcher and distance effects are
constant across all five transitions. Estimating Equation~\ref{eq:period}
instead -- letting both vary freely by period -- and restricting to
the three transitions with enough switch events to identify a
period-specific coefficient (2020--2021 and 2024--2025 have zero
switch events between them after the baseline-rate restriction, so
contribute to the period fixed effects but not to their own
switcher/distance coefficient), gives Table~\ref{tab:period}.

\begin{table}[htbp]
\centering
\caption{Period-specific switcher and distance coefficients}
\label{tab:period}
\begin{tabular}{@{}lcccc@{}}
\toprule
 & \multicolumn{2}{c}{Switcher (binary)} & \multicolumn{2}{c}{Distance (per mile)} \\
Transition & Coef. & $p$ & Coef. & $p$ \\
\midrule
2021--2022 & $-0.648$ & 0.547 & $+0.0014$ & 0.717 \\
2022--2023 & $-1.718$ & 0.449 & $-0.0027$ & 0.802 \\
2023--2024 & $+8.280^{***}$ & 0.008 & $+0.1132^{***}$ & $<0.001$ \\
\bottomrule
\end{tabular}
\vspace{2pt}
\begin{minipage}{0.85\textwidth}
\footnotesize\emph{Notes:} One joint regression, all period
interactions plus period fixed effects and the baseline-rate control;
$N=3{,}794$; SE clustered by state. An F-test rejects that the three
switcher coefficients are equal ($p=0.002$) and that the three
distance coefficients are equal ($p=0.003$): the pooled
Table~\ref{tab:main} coefficient is not simply the same effect
repeated five times. $^{***}p<0.01$.
\end{minipage}
\end{table}

The 2023--2024 transition -- which contributes 75 of the panel's 177
switch events, concentrated in Ohio-border states (Ohio, West
Virginia, Pennsylvania, Kentucky, Indiana) -- has a positive,
significant coefficient on both specifications, the opposite sign from
the pooled result. We investigated this rather than reporting it at
face value, and three things resolve it. First, in raw, unadjusted
terms switcher counties in this transition declined \emph{more} than
controls (mean $\Delta$ overdose rate $-14.5$ vs.\ $-8.5$); the
positive sign appears only after the baseline-rate control, because
switchers in this specific transition have roughly double the 2020
baseline rate of controls (53.4 vs.\ 26.1), a far more severe imbalance
than in other transitions. Second, this period's control pool is
geographically mismatched to its switcher pool -- comparing
Appalachian counties against a national control pool concentrated in
the South Atlantic and West South Central divisions -- and adding
Census-division fixed effects shrinks the switcher coefficient to
$+2.79$ ($p=0.253$) and the distance coefficient to $+0.015$
($p=0.335$), both no longer significant. Third, a handful of
extreme-outlier small counties (implausibly large overdose-rate swings
driven by small population denominators) have outsized leverage on
this single-period fit; excluding the top 5\% by Cook's distance
likewise removes significance ($+1.82$, $p=0.166$). Combining division
fixed effects with winsorizing the baseline-rate control at its 95th
percentile gives $+2.33$ ($p=0.304$) for the switcher coefficient and a
smaller, though still marginally significant, $+0.032$ ($p=0.036$) for
distance. We read this as within-period estimation noise from
isolating a geographically concentrated, high-baseline-rate cohort with
few observations, not as evidence against the pooled result: it draws
on switch events spread across many more states and periods than any
single transition, and Table~\ref{tab:fe}'s state and county
fixed-effects checks already show the continuous specification is not
being driven by any one state or region.

\subsection{Medical-legalization placebo check}

Re-running the identical stacked design with medical- in place of
recreational-marijuana timing, both the switcher ($b=+0.51$,
se$=1.04$, $p=0.62$) and continuous ($b=-0.015$, se$=0.008$, $p=0.08$)
estimates are null -- consistent with the recreational-retail effect
reflecting retail access specifically rather than the design
mechanically detecting any nearby marijuana-policy change
(Figure~\ref{fig:coef}, bottom two rows).

\subsection{Persistence check}

Does the switcher effect continue, flatten, or reverse in the years
after a county crosses the 50-mile threshold? We re-estimate
Equation~\ref{eq:main} with event-time indicators for the switch
transition and each subsequent transition (one, and two-or-more,
transitions after switching), keeping switched counties in the panel
afterward purely for this check (they are otherwise dropped, per
Section~2). With the baseline-rate control, the switch-transition
effect remains significant ($-2.42$, $p=0.001$), but neither
subsequent-transition coefficient is significant (one transition
after: $+1.22$, $p=0.171$; two-plus: $+1.14$, $p=0.157$;
Figure~\ref{fig:persistence}), and a joint test rejects that the lag
effects equal the switch effect ($p=0.006$). The one-year decline
associated with switching does not continue to compound in later
years; if anything, the point estimates hint at partial reversion,
though not significantly so.

\begin{figure}[htbp]
  \centering
  \includegraphics[width=0.85\textwidth]{figures/fig6_persistence.png}
  \caption{Persistence check: switch-transition effect vs. subsequent-transition effects.}
  \label{fig:persistence}
\end{figure}

\subsection{Welfare translation}

National age-specific overdose rates (CDC NCHS Data Query System,
2024) imply a population-weighted average age at overdose death of
46.1 years. We translate the switcher coefficient into years of life
saved using the gap between life expectancy and this average age at
death, rather than the full remaining actuarial life expectancy, since
a person who dies of an overdose at 46 does not have the population's
full remaining life expectancy at stake, only the gap between when
they died and how long they could plausibly have lived absent this
specific cause. Life expectancy is drawn from CDC's state
life-expectancy-at-birth estimates \citep{CDCStateLE2021},
population-weighted across the 177 switcher counties' states (75.9
years), for a gap of 29.7 years per averted death. Applying the
preferred switcher coefficient to the 32.5 million residents of the
177 non-suppressed switcher counties implies approximately 755 averted
deaths (95\% CI: 245 to 1{,}265) over 2020--2025, and roughly 22,500
years of life saved (95\% CI: 7,300 to 37,600;
Figure~\ref{fig:agedist}). This is illustrative of what the point
estimate would mean, not a confirmed welfare effect: it inherits every
caveat on the underlying coefficient, including its greater
sensitivity to fixed-effects controls (Section~\ref{sec:pretrends}).

\begin{figure}[htbp]
  \centering
  \includegraphics[width=\textwidth]{figures/figC_age_distribution.png}
  \caption{National age distribution of drug-overdose deaths, 2024.}
  \label{fig:agedist}
\end{figure}

% ============================================================
% 4. DISCUSSION AND LIMITATIONS
% ============================================================
\section{Discussion and limitations}

Extending the cross-border distance design to a stacked multi-period
panel that uses every state's actual legalization date resolves a
fragility in an earlier two-snapshot version: the continuous per-mile
specification is now sign-consistent with, and more robust than, the
binary switcher result, surviving even county fixed effects. A formal
pre-trends test finds and corrects a mean-reversion problem in the raw
switcher comparison; a persistence check finds the switch-year effect
does not compound in subsequent years; a medical-legalization placebo
check is null, as it should be if the effect is specific to
retail-market access; a comparison against canonical two-way
fixed-effects and Callaway-Sant'Anna estimators shows our stacked
design sitting between a TWFE estimate inflated in the direction its
known bias predicts and a Callaway-Sant'Anna estimate whose own
pre-treatment coefficients independently confirm the same
mean-reversion problem this paper corrects for; and a dispensary-level
robustness check that replaces state-legalization timing with real,
geocoded dispensary opening dates for a 9-state subsample reproduces
the same sign and a comparable magnitude, though with weaker
statistical power reflecting its smaller, partial-coverage sample.

Main limitations, in order of what we judge most consequential: (1)
VSRR data are provisional and small-cell suppression removes more than
half the panel, concentrated where identifying variation is richest;
(2) the 50-mile radius is a design choice, though Appendix~A's
diff-in-discontinuities check finds no sharper discontinuity right at
that threshold; (3) the panel is bounded to 2020--2025 by CDC's
county-level data availability; (4) the welfare translation mixes a
state-level life-expectancy figure with a national age-at-death
figure; (5) the dispensary-level check in Section~\ref{sec:dispensary}
currently covers only 9 of 24 legal states, so should be read as
directionally confirmatory rather than a fully independent
replication; and (6) neither census-tract nor calendar-season fixed
effects are estimable given the county-level, annual (and, for VSRR's
monthly release, 12-month-rolling) structure of the only available
overdose-mortality data. We recommend this be treated as suggestive
evidence of an opioid-substitution channel specific to recreational
retail-market access, meriting confirmation with final,
non-provisional NCHS mortality data and, eventually, full national
dispensary-level coverage.

% ============================================================
% APPENDIX
% ============================================================
\appendix
\renewcommand{\thesection}{A}
\section{Supplementary material}
\label{sec:appendix}

\noindent\emph{Full robustness-check summary.} Table~\ref{tab:robustsummary}
consolidates every fixed-effects and sample-construction check reported
in this paper into one place. Panel A reproduces Table~\ref{tab:fe}
(the checks that bear directly on the pooled, headline specification
in Table~\ref{tab:main}). Panel B reports the additional checks used
only to diagnose and resolve the 2023--2024 period-specific sign
reversal discussed in Section~\ref{sec:periodeffects}; these are
estimated on that single transition in isolation ($n=691$) and were
never substituted into the pooled specification. The dispensary-level
robustness check (Section~\ref{sec:dispensary}, Table~\ref{tab:dispensary})
uses a different underlying sample (real dispensary-opening dates,
9-state coverage) and so is not reproduced in this table, which is
scoped to checks sharing the main $N=3{,}794$ sample.

\begin{table}[htbp]
\centering
\caption{Full robustness-check summary}
\label{tab:robustsummary}
\begin{tabular}{@{}lcccc@{}}
\toprule
 & \multicolumn{2}{c}{Switcher (binary)} & \multicolumn{2}{c}{Distance (per mile)} \\
Specification & Coef. & $p$ & Coef. & $p$ \\
\midrule
\multicolumn{5}{@{}l}{\emph{Panel A: Pooled specification, all periods ($N=3{,}794$) -- bears on Table~\ref{tab:main}}} \\
No controls & $-3.151$ & 0.005 & $+0.0128^{***}$ & $<0.001$ \\
+ Baseline overdose-rate level (preferred) & $-2.320^{***}$ & 0.004 & $+0.0122^{***}$ & $<0.001$ \\
+ State fixed effects & $-1.411$ & 0.144 & $+0.0129^{***}$ & $<0.001$ \\
+ County fixed effects & $-0.999$ & 0.532 & $+0.0122^{**}$ & 0.023 \\
\addlinespace
\multicolumn{5}{@{}l}{\emph{Panel B: 2023--2024 transition only ($n=691$) -- diagnoses Section~\ref{sec:periodeffects}'s sign reversal}} \\
Baseline-rate control only & $+8.280^{***}$ & 0.008 & $+0.1132^{***}$ & $<0.001$ \\
+ Census-division fixed effects & $+2.786$ & 0.253 & $+0.0149$ & 0.335 \\
+ Winsorized baseline rate (95th pctile) & $+4.574^{**}$ & 0.032 & --- & --- \\
+ Division FE \& winsorized baseline & $+2.328$ & 0.304 & $+0.0320^{**}$ & 0.036 \\
Excluding top 5\% by Cook's distance & $+1.822$ & 0.166 & --- & --- \\
\bottomrule
\end{tabular}
\vspace{2pt}
\begin{minipage}{0.95\textwidth}
\footnotesize\emph{Notes:} SE clustered by state throughout. Panel A:
each row a separate regression on the identical $N=3{,}794$ pooled
sample. Panel B: each row a separate regression restricted to the
2023--2024 transition only; ``---'' denotes specifications where the
distance coefficient was not separately estimated. Panel B's raw,
unadjusted group means (no regression) show switcher counties declined
\emph{more} than controls in this transition ($-14.5$ vs.\ $-8.5$
deaths per 100,000); the positive sign in Panel B's first row appears
only once a linear baseline-rate control is applied to a period with
an unusually large baseline-rate imbalance between switchers and
controls (Section~\ref{sec:periodeffects}). $^{**}p<0.05$,
$^{***}p<0.01$.
\end{minipage}
\end{table}

\noindent\emph{Nonlinearity check.} Table~\ref{tab:main}'s continuous
specification imposes a constant marginal effect. Adding a quadratic
term in $\Delta\text{Dist}_{it}$, the quadratic term is itself
significant ($p=0.005$; F-test, $p=0.009$; joint F-test with a cubic
term, $p<0.001$) -- unlike an earlier two-snapshot version of this
design, where no nonlinearity was detectable.
Figure~\ref{fig:doseresponse} plots the dose-response relationship and
implied marginal effect.

\begin{figure}[htbp]
  \centering
  \includegraphics[width=\textwidth]{figures/figA_doseresponse_multiperiod.png}
  \caption{Dose-response (panel a) and implied marginal effect (panel b), pooled stacked-transition panel.}
  \label{fig:doseresponse}
\end{figure}

\noindent\emph{Diff-in-discontinuities check.} We also tested for a
discontinuous jump in overdose mortality exactly at the 50-mile
threshold, with the bandwidth on each side chosen by the CCT
MSE-optimal selector \citep{CalonicoCattaneoTitiunik2014} in each
year's cross-section separately, differencing the 2025 and 2020
estimates following \citet{GrembiNanniciniTroiano2016}. This check is
inconclusive: the only significant year is 2020 (the baseline year,
before almost any rollout had occurred, so not itself interpretable as
a treatment effect), and the discontinuity shrinks rather than
strengthens by 2025 (diff-in-disc $=-14.68$, se$=8.79$, $p=0.095$). We
read this as a reason for some caution about the near/far comparison
generally, not as confirming or overturning the main results.

\noindent\emph{State-level trends and legalization timing.}
Figure~\ref{fig:statetrends} plots monthly overdose mortality,
2015--2026, for all 24 states in the source timing table, with both
the policy-effective date and the first-retail-sales date marked.
Figure~\ref{fig:eventstudy} pools the 23 states with a retail date,
aligned by months since first retail sales; neither figure feeds into
the county-level regression design.

\begin{figure}[htbp]
  \centering
  \includegraphics[width=\textwidth]{figures/figD_state_overdose_trends_new.png}
  \caption{State overdose mortality trends, 2015--2026, with policy-effective and retail-sales dates.}
  \label{fig:statetrends}
\end{figure}

\begin{figure}[htbp]
  \centering
  \includegraphics[width=\textwidth]{figures/figE_eventstudy_overdose_new.png}
  \caption{Overdose mortality around the retail-legalization cutoff, pooled across 23 states.}
  \label{fig:eventstudy}
\end{figure}

% ============================================================
% REFERENCES
% ============================================================
\bibliographystyle{apalike}
\bibliography{references}

\end{document}
